% Part: conditional-logics
% Chapter: minimal-change-semantics
% Section: sphere-models

\documentclass[../../../include/open-logic-section]{subfiles}

\begin{document}

\olfileid[pt]{con}{min}{sph}

\olsection{Modelos de esferas}

Uma maneira de fornecer uma semântica formal para os contrafactuais é transformar
a explicação informal de Lewis em uma estrutura matemática. As esferas
ao redor de um mundo~$w$ são, então, conjuntos de mundos. Como as esferas são
aninhadas, os conjuntos de mundos ao redor de~$w$ precisam estar linearmente ordenados pela
relação de inclusão.

\begin{defn}
  Um \emph{modelo de esferas} é uma tripla~$\mModel{M} = \tuple{W, O, V}$,
  em que $W$ é um conjunto não vazio de mundos, $V\colon \PVar \to \Pow{W}$
  é uma valoração, e $O\colon W \to \Pow{\Pow{W}}$ atribui a cada
  mundo~$w$ um \emph{sistema de esferas}~$O_w$. Para cada $w$, $O_w$ é um
  conjunto de conjuntos de mundos e deve satisfazer:
  \begin{enumerate}
  \item $O_w$ é \emph{centrado} em~$w$: $\{w\} \in O_w$.
  \item $O_w$ é \emph{aninhado}: sempre que $S_1$, $S_2 \in O_w$, $S_1
    \subseteq S_2$ ou $S_2 \subseteq S_1$, isto é, $O_w$ é linearmente
    ordenado por~$\subseteq$.
  \item $O_w$ é fechado sob uniões não vazias.
  \item $O_w$ é fechado sob interseções não vazias.
  \end{enumerate}
\end{defn}

A intuição por trás de $O_w$ é que os mundos ``ao redor'' de $w$ são
estratificados de acordo com sua distância de~$w$. A esfera mais interna
é apenas o próprio $w$, isto é, o conjunto~$\{w\}$: $w$ está mais próximo
de~$w$ do que os mundos em qualquer outra esfera. Se $S \subsetneq S'$, então
os mundos em $S' \setminus S$ estão mais distantes de $w$ do que os
mundos em~$S$: $S' \setminus S$ é a ``camada'' entre $S$ e
os mundos fora de~$S'$. Em particular, devemos conceber as
esferas como contendo todos os mundos dentro de sua superfície externa;
elas não são apenas as camadas individuais.

\begin{figure}
\begin{center}
\begin{tikzpicture}[layerwidth=1.5,scale=.6]\tiny
  \spheresystem{3}
  \propositionintersect{0}{3}{60}{6}
  \path (0,0) node[world] {$w$};
  \spherepos{-30}{2}{node[world] {$w_2$}}
  \spherepos{220}{2}{node[world] {$w_3$}}
  \spherepos{90}{2}{node[world] {$w_1$}}
  \spherepos[shift={(.1,0)}]{10}{3}{node[world] {$w_5$}}
  \spherepos[shift={(.1,0)}]{-10}{3}{node[world] {$w_6$}}
  \spherepos{225}{3}{node[world] {$w_4$}}
  \spherepos{20}{4}{node[world] {$w_7$}}
  \path (5,0) node[left] {$p$};
\end{tikzpicture}
\caption{Diagrama de um modelo de esferas}
\ollabel{fig:sphere-model}
\end{center}
\end{figure}

O diagrama em \olref{fig:sphere-model} corresponde ao modelo de
esferas com $W = \{w, w_1, \dots, w_7\}$, $V(p) = \{w_5, w_6,
w_7\}$. A esfera mais interna é $S_1 = \{w\}$. Os mundos mais próximos de $w$
são $w_1, w_2, w_3$, portanto a esfera seguinte, maior, é $S_2 = \{w, w_1,
w_2, w_3\}$. Os mundos mais afastados são $w_4$, $w_5$, $w_6$, de modo que a
esfera mais externa é $S_3 = \{w, w_1, \dots, w_6\}$. O sistema de
esferas ao redor de $w$ é $O_w = \{S_1, S_2, S_3\}$. O mundo~$w_7$ não está
em nenhuma esfera ao redor de~$w$. Os mundos mais próximos nos quais $p$~é verdadeiro
são $w_5$ e~$w_6$; assim, a menor esfera que admite $p$
é~$S_3$.

Para definir a satisfação de uma fórmula~$!A$ em um mundo~$w$ de um modelo de
esferas~$\mModel M$, $\mSat{M}{!A}[w]$, ampliamos a definição para
!!{formula}s modais, incluindo uma cláusula para $!B \cif !C$:

\begin{defn}
  $\mSat{M}{!B \cif !C}[w]$ se, e somente se, ou
  \begin{enumerate}
  \item\ollabel{sphere-vac} para todo $u \in \bigcup O_w$, $\mSat/{M}{!B}[u]$; ou
  \item\ollabel{sphere-nonvac} para algum $S \in O_w$,
    \begin{enumerate}
      \item $\mSat{M}{!B}[u]$ para algum $u \in
        S$; e
      \item para todo $v \in S$, ou
        $\mSat/{M}{!B}[v]$ ou $\mSat{M}{!C}[v]$.
    \end{enumerate}
  \end{enumerate}
\end{defn}

Segundo essa definição, $\mSat{M}{!B \cif !C}[w]$ se, e somente se, ou o
antecedente~$!B$ é falso em todos os pontos das esferas ao redor de~$w$, ou
existe uma esfera~$S$ na qual $!B$ é verdadeiro e o condicional material
$!B \lif !C$ é verdadeiro em todos os mundos dessa esfera que
``admite $!B$''. Observe que não exigimos na definição que $S$ seja a
esfera \emph{mais interna} que admite $!B$, ao contrário do que se poderia
esperar da explicação intuitiva. Mas, se a condição
em~\olref{sphere-nonvac} é satisfeita por alguma esfera~$S$, então ela também é
satisfeita por todas as esferas contidas em~$S$ e, portanto, em particular,
pela esfera mais interna.

Observe também que a definição de modelos de esferas não exige que
\emph{exista} uma esfera mais interna que admita $!B$: podemos ter uma
sequência infinita $S_1 \supsetneq S_2 \supsetneq \dots \supsetneq
\{w\}$ de esferas que admitem $!B$ e, portanto, nenhuma esfera mais interna
que admita $!B$. Nesse caso, $\mSat{M}{!B \cif !C}[w]$ se, e somente se,
$!B \lif !C$ vale em todas as esferas $S_i$, $S_{i+1}$, \dots, para
algum~$i$.

\end{document}
